\documentclass[10pt,a4paper]{amsart}
\usepackage[margin=19mm]{geometry}
\usepackage{amsmath,amssymb,mathtools}
\usepackage[scr=rsfs,cal=euler]{mathalfa}
\usepackage{xfrac}
\usepackage[unicode,hidelinks,bookmarks=false]{hyperref}
\newcommand{\ie}{i.e.\ }
\newcommand{\cf}{cf.\ }
\newcommand{\resp}{respectively\ }
\newcommand{\Subsection}{\subsection}
\providecommand{\nobreakdash}{\nobreak\hbox{-}}
\setlength{\emergencystretch}{2em}
\multlinegap0pt
\usepackage{amssymb}
\usepackage{mathtools}
\usepackage[shortlabels]{enumitem}
\usepackage{xcolor}
\newtheorem{lemma}{Lemma}
\newtheorem{theorem}{Theorem}
\title{\bfseries A Nonlinear Deficiency Identity for the Riemann Zeta Function with Optimal Approximation Rates}
\author{Meisam Mohammady}
\date{Source: arXiv:2604.16530v3 (CC BY 4.0)}
\begin{document}
\maketitle

\noindent\emph{Source and licence.} \bfseries A Nonlinear Deficiency Identity for the Riemann Zeta Function with Optimal Approximation Rates. Version arXiv:2604.16530v3 (CC BY 4.0), distributed by arXiv under the Creative Commons Attribution 4.0 International licence, http://creativecommons.org/licenses/by/4.0/, as recorded in the arXiv licence metadata for this exact version and shown on the abstract page https://arxiv.org/abs/2604.16530v3. Full licence notice: \texttt{licenses/arxiv-2604.16530-CC-BY-4.0.txt}.\par\smallskip

\noindent\textit{Editorial scope.} Counted unit: the article's theorem thm:general (source0.tex lines 1078--1096). The article's complete original proof of it is delivered with it (source0.tex lines 1098--1142, one \texttt{proof} environment of 45 lines). No mathematics is written, repaired or re-derived: the statement and the proof are the article's own lines, in the article's own notation, and no retained line is substituted or re-flowed.  Retained uncounted as the source-local context the proof needs: lines 594--602; lines 1057--1068.  Nothing was omitted from any retained span, so the package carries no omission record.  No retained line contains a citation, so the wrapper carries no bibliography and no bibliography entry.  No retained line contains a figure environment, an image-inclusion command or drawing code, so no image is delivered and the package declares no asset dependency.  Editorial formatting only, in addition: an AMS \texttt{amsart} preamble that restates the article's own declarations and macros used by the retained lines, namely the environments \texttt{lemma}, \texttt{theorem} on separate counters, together with the standard packages those lines need.  The retained environments print in this document as Lemma 1, Theorem 1 in order of appearance, whereas the article prints the same items with the numbering of its own full document.  The complete native source of the article as distributed in the arXiv e-print for this exact version is bundled unchanged as \texttt{sources/198679/source0.tex}.
\begin{lemma}[Tail Asymptotics]
\label{lem:tail}
For every \(p>1\),
\[
t_n^{(p)}
=
\frac{1}{(p-1)n^{p-1}}+O(n^{-p}).
\]
\end{lemma}

\begin{lemma}
\label{lem:taylor}
Let $a>1$ and $t_n \to 0$. Then
\[
(\zeta(p) - t_n)^a
=
\zeta(p)^a
- a\zeta(p)^{a-1} t_n
+ \frac{a(a-1)}{2}\zeta(p)^{a-2} t_n^2
+ O(t_n^3).
\]
\end{lemma}

\begin{theorem}
\label{thm:general}
Let $p>1$ and $q>p$. Define
\[
B_n^{(p,q)}
=
\zeta(p)^{q/p}
-
D_n^{(p,q)}
-
\frac{q}{p}\zeta(p)^{q/p-1}(\zeta(p)-S_n^{(p)}).
\]
Then
\[
B_n^{(p,q)} - \zeta(q)
=
O\!\left(n^{-\min(2p-2,\;q-1)}\right).
\]
\end{theorem}

\begin{proof}
Let
\[
t_n^{(p)} = \zeta(p) - S_n^{(p)}.
\]

Using Lemma~\ref{lem:taylor} with $a=q/p$,
\[
\zeta(p)^{q/p} - S_n^{(p)\,q/p}
=
\frac{q}{p}\zeta(p)^{q/p-1} t_n^{(p)}
-
\frac{q(q-p)}{2p^2}\zeta(p)^{q/p-2} t_n^{(p)\,2}
+ O(t_n^{(p)\,3}).
\]

Substituting into the estimator,
\[
B_n^{(p,q)} - \zeta(q)
=
-\frac{q(q-p)}{2p^2}\zeta(p)^{q/p-2} t_n^{(p)\,2}
- R_n^{(q)} + O(t_n^{(p)\,3}),
\]
where
\[
R_n^{(q)} = \zeta(q) - \sum_{k=1}^n \frac{1}{k^q}.
\]

By Lemma~\ref{lem:tail},
\[
t_n^{(p)} \sim n^{-(p-1)}, \quad
R_n^{(q)} \sim n^{-(q-1)}.
\]

Thus
\[
t_n^{(p)\,2} \sim n^{-(2p-2)}.
\]

Therefore the dominant term is
\[
\min(2p-2,\;q-1),
\]
which proves the result.
\end{proof}

\end{document}
